\documentclass[12pt]{article}
\usepackage[utf8]{inputenc}
\usepackage[T1]{fontenc}
\usepackage[margin=1in]{geometry}
\usepackage{setspace}
\usepackage{xcolor}
\usepackage[colorlinks=true, urlcolor=blue!60!black, linkcolor=blue!60!black]{hyperref}
\setlength{\parindent}{0em}
\setlength{\parskip}{0.6em}
\renewcommand{\baselinestretch}{1.0}
\begin{document}
\begin{center}
{\LARGE \textit{Penn Central Transportation Co.\ v.\ New York City}, 438 U.S.\ 104 (1978)\par}
\vspace{0.5em}
{\large Case Brief --- Two-Page Edition\par}
\end{center}

\section{Memory Jogger}
No taking when New York barred a tower above Grand Central Terminal; regulatory takings turn on an ad hoc weighing of three factors against the parcel as a whole (\href{https://supreme.justia.com/cases/federal/us/438/104/}{Justia}).

\section{Detailed Case Facts}
After New York designated Grand Central Terminal a landmark, its Landmarks Preservation Commission rejected Penn Central's plans for a 53- or 55-story office tower above it, calling one an ``aesthetic joke.'' Landmark owners could transfer unused air rights (TDRs) to nearby lots.

\section{Procedural History}
The trial court granted relief; the Appellate Division reversed, and the New York Court of Appeals affirmed.

\section{Judicial Votes}
Affirmed, 6--3. Brennan wrote for Stewart, White, Marshall, Blackmun, and Powell; Rehnquist dissented, joined by Burger and Stevens.

\section{Holding}
No taking. Courts weigh economic impact, interference with distinct investment-backed expectations, and the character of the government action, measured against the parcel as a whole rather than severed air rights.

\section{Analysis of Opinions}
\textbf{Majority:} There is no ``set formula''; physical invasions are more readily takings than programs adjusting economic burdens. Penn Central could not sever its air rights; the citywide landmark program was not spot zoning; and the Terminal kept its existing use, a reasonable return, the chance of a smaller addition, and valuable TDRs. \textbf{Dissent:} unlike zoning's ``average reciprocity of advantage,'' the law singled out under 0.1\% of buildings because ``too good a job was done'' building them.

\section{Examples: Future Applications}
\textbf{Same-Side:} denying a tower atop a profitable landmark warehouse with TDRs. \textbf{Opposite-Side:} a public easement over a private dock, a \textit{per se} taking under \textit{Cedar Point}. \textbf{Fence-Sitter:} landmarking one small building, cutting its value 75 percent, with no TDRs.

\section{Critique}
The unweighted factors are unpredictable and favor regulators, and the ``parcel as a whole'' denominator remains fluid.

\section{Key Quotations}
\begin{quote}
``In engaging in these essentially ad hoc, factual inquiries, the Court's decisions have identified several factors that have particular significance. The economic impact of the regulation on the claimant and, particularly, the extent to which the regulation has interfered with distinct investment-backed expectations are, of course, relevant considerations. \dots\ So, too, is the character of the governmental action. A `taking' may more readily be found when the interference with property can be characterized as a physical invasion by government \dots\ than when interference arises from some public program adjusting the benefits and burdens of economic life to promote the common good.''
\end{quote}
\noindent --- Justice Brennan, for the Court, 438 U.S.\ at 124.

\end{document}
